% 発表原稿（日本語・簡易版） — 1スライドあたり2行程度
% ビルド: lualatex speaker_script_ja.tex
\documentclass[11pt,a4paper]{ltjsarticle}
\usepackage[a4paper,margin=2.2cm]{geometry}
\usepackage{luatexja-fontspec}
% IPAGothic は単一ウェイトなので、太字は疑似ボールドで出す
\setmainjfont{IPAGothic}[AutoFakeBold=2.5]
\setsansjfont{IPAGothic}[AutoFakeBold=2.5]
\usepackage{amsmath}
\usepackage{xcolor}
\usepackage{tcolorbox}
\usepackage{fancyhdr}
\usepackage{parskip}
\usepackage{setspace}

\definecolor{luhblue}{RGB}{21,101,192}
\definecolor{alertred}{RGB}{198,40,40}
\definecolor{soft}{RGB}{245,247,250}

\setstretch{1.15}
\pagestyle{fancy}\fancyhf{}
\lhead{\small 確率有限要素法 口頭発表 --- 日本語原稿（簡易版）}
\rhead{\small\thepage}
\renewcommand{\headrulewidth}{0.4pt}

% \slide{番号}{タイトル}{時間}
\newcommand{\slide}[3]{%
  \medskip
  \noindent\colorbox{luhblue}{\parbox{\dimexpr\linewidth-2\fboxsep}{%
    \color{white}\bfseries\small スライド #1 \quad|\quad #2 \hfill #3}}\par
  \nopagebreak\vspace{0.25em}}

\begin{document}

\begin{center}
{\LARGE\bfseries 発表原稿（日本語・簡易版）}\\[0.4em]
{\large 製造ばらつきを考慮した JAXA H3 フェアリングの信頼性評価}\\[0.3em]
{\normalsize 西岡 佳祐 \quad 学籍番号 10081049 \quad LUH 2026年夏学期}
\end{center}

\vspace{0.4em}
\begin{tcolorbox}[colback=soft,colframe=luhblue,boxrule=0.6pt,arc=2pt]
{\bfseries この原稿について}\\[0.3em]
英語スライド18枚（表紙＋17枚）に対応した、1枚あたり2行程度の日本語要点集です。
英語の読み上げ原稿 \texttt{speaker\_script\_oral.pdf} を逐語訳したものではなく、
各スライドで何を言うべきかを圧縮したものです。時間表記は英語版スライドの進行に合わせてあります。
\end{tcolorbox}

%----------------------------------------------------------
\slide{1}{表紙}{0:00 -- 0:20}

西岡佳祐です。確率有限要素法の最終プロジェクトを発表します。

問いは単純で、複合材フェアリングの材料が製造でばらつくとき、それが構造応答にどう伝播し、構造は信頼できるのか、ということです。

%----------------------------------------------------------
\slide{2}{背景：複合材フェアリングの製造ばらつき}{0:20 -- 1:05}

CFRP 表皮と接着層の物性は、繊維体積率・硬化条件・接着厚のばらつきで本質的に変動します。この変動が応力・損傷・変形、ひいては信頼性に伝播します。

公称値だけの決定論的解析では、計算された余裕のうちどれだけが実在のもので、どれだけが仮定した数値の産物なのかが分かりません。

%----------------------------------------------------------
\slide{3}{対象構造：H3 ペイロードフェアリング}{1:05 -- 1:45}

対象は CFRP 表皮とアルミハニカムコアのサンドイッチパネルです。直径 5.2 メートル、表皮 0.6 から 1.4 ミリ、コア 8 ミリ。

ばらつきが宿るのは表皮の剛性と接着層で、これが5つの不確実パラメータになります。接着界面は凝集力モデルで表現しています。

%----------------------------------------------------------
\slide{4}{全体の流れ}{1:45 -- 2:25}

左が5つの材料パラメータで、それぞれ分布として与えます。中央が決定論的な有限要素モデルで、これを選んだ設計点で評価します。

右が取り出すもので、ピーク応力の確率分布と Sobol 感度指数の2つです。以降はこの3段階の詳細になります。

%----------------------------------------------------------
\slide{5}{手法①：多項式カオス展開}{2:25 -- 3:25}

応答を直交多項式の有限和で近似します。ガウス入力に対してはエルミート多項式で、基底は入力の測度で決まり、任意に選ぶものではありません。

重要なのは多項式の形ではなく直交性です。直交しているおかげで、平均は0次の係数、分散は残りの係数の二乗和、Sobol 指数もその部分和として、追加の解析なしに得られます。

%----------------------------------------------------------
\slide{6}{手法②：Smolyak 疎求積}{3:25 -- 4:15}

全テンソル積は $n^d$ 点必要で、5次元では現実的ではありません。Smolyak は総次数空間でのみ厳密性を保つことで、2次66点、3次286点に削減します。

実装上の注意が一つあります。疎な規則は一般に負の重みを生むので、係数は求積射影ではなく最小二乗回帰で当てています。66点に対し未知数21なので優決定系で、Leave-one-out 検証も可能です。

%----------------------------------------------------------
\slide{7}{5つの不確実パラメータ}{4:15 -- 5:00}

表皮に2つ、繊維方向弾性率 $E_1$ が146ギガパスカルで変動係数5パーセント、せん断弾性率 $G_{12}$ が10パーセント。

凝集力則に3つ、法線剛性・Mode-I 破壊エネルギー・強度が15から20パーセント。剛性と破壊エネルギーは正値なので、公称値の $\pm50$ パーセントで切断した正規分布としています。

%----------------------------------------------------------
\slide{8}{FE モデルと注目量}{5:00 -- 5:45}

Abaqus/Standard による静的2ステップ解析で、Step-1 が圧力、Step-2 が増分荷重です。基部は完全拘束しています。

注目量は3つ、表皮の最大 von Mises 応力（限界1200メガパスカル）、最大凝集損傷 SDEG、最大変位です。各設計点で入力ファイルの材料カードだけを書き換えるので、完全に非介入型です。

%----------------------------------------------------------
\slide{9}{1回の解析はこう見える}{5:45 -- 6:20}

1回の評価がどういうものかを示したのがこの応力場です。求積点1点が、この決定論的解析1回に対応します。

応力はアクセスドア開口部まわりに集中します。全体モデルは文脈のために示しており、確率解析自体は代表パネルを対象にしています。

%----------------------------------------------------------
\slide{10}{結果①：2次で統計量が収束}{6:20 -- 7:05}

ピーク応力は2次でも3次でも209.7メガパスカル、標準偏差は2.375と2.378で、有効数字4桁まで一致しました。

これは単なる収束確認以上の意味を持ちます。高次項に依存していれば追加の220回が標準偏差を動かしたはずで、動かなかった。つまり2次は節約策ではなく、この応答にとって正しいモデル次数だということです。

%----------------------------------------------------------
\slide{11}{結果②：$E_1$ が分散の99\%超を支配}{7:05 -- 7:55}

繊維方向弾性率だけで応力分散の99.2パーセント、変位分散の99.95パーセントを説明します。他のどのパラメータも1パーセントに届きません。

$G_{12}$ は変動係数が $E_1$ の2倍あるのに寄与は1パーセント未満です。また1次と全次の指数がほぼ等しく、入力が加法的に働き、検出できる交互作用がないことを意味します。

%----------------------------------------------------------
\slide{12}{なぜ $E_1$ が支配するのか}{7:55 -- 8:50}

1次近似では、出力分散は感度の二乗と入力分散の積の和です。各パラメータは感度とばらつきの積で効くのであって、どちらか一方では効きません。

パネルは膜応力で荷重を受けるため、応力も変位も軸方向剛性で決まり、$E_1$ の感度が大きくなります。$G_{12}$ はその反例で、ばらつきは2倍でも感度が小さいため出力に現れません。順位を決めるのは、最もばらつく量ではなく、構造が増幅するばらつきです。

%----------------------------------------------------------
\slide{13}{結果③：PCE とニューラルネットの比較}{8:50 -- 9:45}

66点では両者とも $10^{-3}$ メガパスカル程度で、この指標では NN の方がわずかに良いくらいです。つまり設計点での精度は決め手になりません。安定性が決め手です。

訓練点を66から286に増やすと NN の誤差は25倍に増えます。数千の重みを持つネットワークは、写像がこれほど滑らかだと過学習します。PCE 2次は維持、3次はわずかに悪化で、これは軽度の過剰パラメータ化です。

%----------------------------------------------------------
\slide{14}{信頼性：広い余裕と、同定できない界面}{9:45 -- 10:40}

平均応力は閾値の416標準偏差下にあり、破壊確率はゼロ、信頼性指標は無限大です。ただしゼロは保証ではなく、解析の分解能以下と読むべきです。

もう一つ帰結があります。損傷が全352回でゼロだったため、凝集力の3パラメータは構造的に同定不能です。出力を動かさないので、どれだけサンプルを増やしても指数はゼロのままで、5次元問題はこの荷重では実質1次元に潰れています。

%----------------------------------------------------------
\slide{15}{発展課題 T9：非ガウス KL 確率場}{10:40 -- 11:50}

$E_1$ をスカラーではなく確率場とし、異方性指数カーネルで Karhunen--Lo\`{e}ve 展開しました。ガウス場は負の剛性を許すので、無記憶変換で対数正規の変換場にし、正値性を構成的に保証します。

変換は順位相関を保ちますが線形相関は保ちません。$E_1$ の5パーセントのばらつきでは歪みは $3\times10^{-4}$ で、ここでの実益は正値性です。空間平均により実効剛性のばらつきは7300から約4700メガパスカルに下がりますが、これは入力側の結果です。ピーク応力は局所極値なので、この保守性はそのままは転用できません。

%----------------------------------------------------------
\slide{16}{発展課題 T13：正しい事前知識が柔軟性に勝る}{11:50 -- 12:55}

PCE・ガウス過程・ニューラルネットを10点から286点で比較しました。ただし参照応答自体が2次多項式なので、2次展開はそれを厳密に含みます。

$10^{-13}$ という値は厳密表現可能性を測っており、汎化誤差ではありません。意味があるのはガウス過程と NN の約500倍の差です。測っているのは、データが乏しいときに正しい事前知識が持つ価値であって、手法の普遍的な優劣ではありません。

%----------------------------------------------------------
\slide{17}{発展課題 T6：ベイズ同定}{12:55 -- 13:55}

$E_1$ が分散の99パーセントを説明するので、5つの中で応答計測から最も同定しやすいのも $E_1$ です。PCE の線形項を尤度とする共役ガウス問題として解くと、事後分布は閉形式で得られます。

事後標準偏差は7300から、5点で2582、10点で1885（74パーセント減）、20点で1356まで下がります。ただし減るのは認識論的不確実性であり、バッチ間の偶然的ばらつきは不変です。破壊確率は既にゼロなので信頼性評価は改善しません。価値は公差配分の側にあります。

%----------------------------------------------------------
\slide{18}{結論と限界}{13:55 -- 15:00}

66回の解析による2次展開で応答統計量は捉えられ、これは節約ではなく正しいモデル次数だと考えます。実務的な結論は、繊維方向弾性率が構造ばらつきの99パーセント超を支配するということで、品質管理をここに集中させるのが最も効果的です。

限界を述べて終わります。破壊確率ゼロは分解能の限界、入力は独立と仮定、パネルは1枚のみ、そして確率場の実現値をソルバーに流していません。自然な続きは、凝集力域に入る荷重レベルで、相関を入れた入力モデルにより再実施することです。

ご清聴ありがとうございました。質問をお受けします。

%----------------------------------------------------------
\vspace{1em}
\begin{tcolorbox}[colback=soft,colframe=alertred,boxrule=0.6pt,arc=2pt]
{\bfseries 突かれやすい3点}\\[0.4em]
\textbf{荷重が低すぎるのでは} --- その通りです。設計荷重では凝集力域に入らず、接着層のパラメータは同定不能です。これは本研究の最大の限界で、損傷開始まで荷重を上げるか、欠陥・衝撃シナリオを入れる必要があります。\\[0.4em]
\textbf{確率場は FE に入れたのか} --- 入れていません。Topic 9 はスタンドアロンの数値実験で、実現値をソルバーに流す検証は未実施です。\\[0.4em]
\textbf{入力は本当に独立か} --- 独立と仮定しましたが、$E_1$ と $G_{12}$ は同じ積層板の物性なので実際には相関があるはずです。相関を入れた入力モデルは今後の課題です。
\end{tcolorbox}

\end{document}
